\ProvidesFile{Dependencies.tex}[Зависимости]


\subsection{Зависимости объектов и функций}

Так часто бывает что при исполнении какой-либо функции или метода объекта, ему нужно вызвать методы другого объекта, который решает вспомогательную задачу.
Например, метод создающий матрицу может использовать случайный генератор для инициализации ее элементов.
Или класс для отображения пользовательского интерфейса может использовать класс-локализатор, который знает все надписи для отображения.
Давайте обсудим, каким образом можно передавать зависимости в функции и в объекты классов.

\subsubsection{Зависимости функций}

Начнем с простого примера функции:
\begin{cppcode}
Matrix create_matrix(Rows m, Cols n);
\end{cppcode}
Мы подразумеваем, что этот метод создает случайную матрицу указанных размеров.%
\footnote{Здесь \verb"Rows" и \verb"Cols" -- strict aliases для размеров матрицы, подробнее в разделе~\ref{section::CtorArgs}.}
Теперь вопрос в том, какой случайный генератор будет использовать функция и как это сделать правильно?

Что мы можем сделать:
\begin{enumerate}
\item Передать случайный генератор явно в виде аргументов.
При этом у функции технически два вида аргументов:
\begin{enumerate}
\item Шаблонные параметры.
\begin{cppcode}
class Random { ... };

template<class Rnd>
Matrix create_matrix(Rows m, Cols n);

Matrix m = create_matrix<Random>(3, 4);
\end{cppcode}
Эти аргументы должны быть известны на этапе компиляции и не меняются во время исполнения.
Так же обратим внимание, что тут передается не конкретный объект, который помнит свое состояние, а имя класса.
Если сделать наивную имплементацию вида
\begin{cppcode}
template<class Rnd>
Matrix create_matrix(Rows m, Cols n) {
  Rnd rnd{123};
  Matrix mat(m, n);
  for (auto& element : mat.data())
    element = rnd.uniform(0., 1.);
  return mat;
}
\end{cppcode}
То мы теряем всякую случайность.
Потому что при каждом вызове метода мы создаем полностью предетерминированный случайный генератор, инициализированный одинаково, и потому инициализирующий все матрицы одинаково.
В таком случае у нас должен быть статический метод, который помнит предыдущее состояние:
\begin{cppcode}
template<class Rnd>
Matrix create_matrix(Rows m, Cols n) {
  Matrix mat(m, n);
  for (auto& element : mat.data())
    element = Rnd::uniform(0., 1.);
  return mat;
}
\end{cppcode}
Метод \verb"uniform" тогда обязан хранить некоторое состояние внутри \verb"static" переменной, то есть глобальной переменной, которая видна только этому методу:
\begin{cppcode}
class Random {
public:
  ...
  static Matrix uniform(double a, double b) {
    static Generator generator{123};
    return uniform_impl(a, b, generator);
  }
};
\end{cppcode}

\item Аргументы функции.
\begin{cppcode}
class Random { ... };

Matrix create_matrix(Rows m, Cols n, Random& rnd);

Random rnd{123};
Matrix m = create_matrix(3, 4, rnd);
\end{cppcode}
Эти аргументы могут меняться во время исполнения.
Это по сути единственное отличие от предыдущего случая. 
\end{enumerate}

\item Использовать глобальный случайный генератор.
И тут есть две версии:
\begin{enumerate}
\item Использовать приватный \verb"static" генератор.
\begin{cppcode}
Matrix create_matrix(Rows m, Cols n) {
  static Random rnd{123};
  Matrix mat(m, n);
  for (auto& element : mat.data())
    element = rnd.uniform(0., 1.);
  return mat;
}
\end{cppcode}

\item Использовать Singleton.
\begin{cppcode}
Matrix create_matrix(Rows m, Cols n) {
  GRndAccessor rnd;
  Matrix mat(m, n);
  for (auto& element : mat.data())
    element = rnd->uniform(0., 1.);
  return mat;
}
\end{cppcode}
Подробнее о том, что тут происходит можно посмотреть в разделе~\ref{section::Global}, что такое \verb"Accessor" можно глянуть в подразделе~\ref{section::SingletonParams}.
Разница между первым и вторым подходом только в одном: в первом случае случайный генератор виден только одной функции, а во втором случае несколько функций могут обращаться к одному и тому же генератору.
\end{enumerate}
\end{enumerate}

Вот так выглядят вызовы во всех $4$ случаях:
\begin{cppcode}
Matrix m = create_matrix<Random>(3, 4); // explicit compile-time
Matrix m = create_matrix(3, 4, rnd);    // explicit run-time
Matrix m = create_matrix(3, 4);         // implicit static
Matrix m = create_matrix(3, 4);         // implicit singleton
\end{cppcode}

Сделаем несколько замечаний
\begin{enumerate}
\item В случаях $1$ и $2$ нам приходится явно указывать генератор, что может быть не удобно.

\item В случаях $3$ и $4$ нет никакой возможности понять вообще используется ли какой-то внешний  генератор или нет.
И эти два случая не отличимы на стороне пользователя.

\item В случаях $1$ и $2$ можно подменять генераторы при разных вызовах метода.
\begin{cppcode}
Matrix m = create_matrix<Random1>(3, 4); // explicit compile-time
Matrix m = create_matrix<Random2>(3, 4); // explicit compile-time
Matrix m = create_matrix(3, 4, rnd1);    // explicit run-time
Matrix m = create_matrix(3, 4, rnd2);    // explicit run-time
\end{cppcode}

\item В случаях $3$ и $4$ не возможно подменить один генератор другим, при всех вызовах используется один и тот же генератор.

\item В случаях $1$ и $2$ можно избавиться от необходимости вводить явно генератор на стороне пользователя, указав дефолтный генератор.
Для шаблонного случая вот так:
\begin{cppcode}
template<class Rnd>
Matrix create_matrix(Rows m, Cols n);  // template function

Matrix create_matrix(Rows m, Cols n) { // overloaded function
  return create_matrix<DefRandom>(m, n);
}

Matrix m = create_matrix<Random>(3, 4);
Matrix m = create_matrix(3, 4);
\end{cppcode}
Либо с помощью дефолтных параметров аргумента:
\begin{cppcode}
template<class Rnd = Random>
Matrix create_matrix(Rows m, Cols n);

Matrix m = create_matrix<Random>(3, 4);
Matrix m = create_matrix(3, 4);
\end{cppcode}
Для аргументов метода вот так:
\begin{cppcode}
Matrix create_matrix(Rows m, Cols n, Random& rnd);

Matrix create_matrix(Rows m, Cols n) {
  return create_matrix(m, n, GRandom::instance());
}

Matrix m = create_matrix(3, 4, rnd);
Matrix m = create_matrix(3, 4);
\end{cppcode}
Либо с помощью дефолтных параметров аргумента:
\begin{cppcode}
Matrix create_matrix(Rows m, Cols n, Random& rnd = GRandom::instance());

Matrix m = create_matrix(3, 4, rnd);
Matrix m = create_matrix(3, 4);
\end{cppcode}
С точки зрения API оба подхода эквивалентны.
Однако, разница появляется на уровне ABI.
Если вам это важно, то лучше пользоваться первым подходом.
Так же этот вопрос подробно обсуждается в разделе~\ref{section::SingletonRefactoring}.
\end{enumerate}
Что из этого использовать?
Я бы использовал подходы $1$ и $2$ с перевесом в сторону $2$, если вам нужно менять что-либо в run-time.
Лучше для функций придумать нельзя, просто потому что для функции есть только два вида внешних данных: ее аргументы и глобальные переменные.
И мы по сути рассмотрели все возможные подходы.
Явная передача данных намного лучше, потому что с ней идет параллельно локальность.
Ситуации, когда надо использовать $3$ и $4$ подход обсуждаются в последних разделах~\ref{section::Global}.
Если вам не достаточно того, что могут себе позволить функции, то вероятно вам лучше инкапсюлировать ее работу в объект класса.
А в этом случае у вас появляются новые способы задания зависимостей, которые будут обсуждаться ниже.
В целом метод класса хорош тем, что у вас кроме глобальных и локальных переменных, появляются переменные класса, которые являются <<контролируемыми глобальными>> ресурсами в рамках одного объекта.

\subsubsection{Зависимости объектов}

Начнем с простого класса:
\begin{cppcode}
class Widget {
public:
  ...
private:
  ...
  Localizator loc_;
};
\end{cppcode}
В данном случае у нас есть класс \verb"Widget", который представляет пользователю какой-то графический интерфейс и за текст в открывшихся окнах отвечает объект \verb"loc_" типа \verb"Localizator".
Мы считаем, что для каждого языка у нас свой локализатор.
А потому смена языка означает смену локализатора.
Теперь обсудим, как можно сообщать эту зависимость объекту:
\begin{enumerate}
\item В конструкторе.
Здесь все просто
\begin{cppcode}
class Widget {
public:
  Widget(const Localizer& loc) : loc_(loc) {}
  ...
private:
  ...
  Localizator loc_;
};
\end{cppcode}
В этом случае есть возможность установить локализатор, но нет возможности его изменить.
Так же нельзя создать \verb"Widget" без локализатора.
Последнее может быть как плюсом, так и минусом.

\item Специальный метод.
\begin{cppcode}
class Widget {
public:
  void set_localizer(const Localizer& loc) {
    loc_ = loc;
  }
  ...
private:
  ...
  std::optional<Localizator> loc_;
};
\end{cppcode}
По сути то же самое, что и в случае конструктора, только теперь мы можем локализатор менять налету.
Тут есть одна тонкость: а чем инициализировать в дефолтном конструкторе?
И тут есть два пути: либо чем-то дефолтным, либо ничем. Но тогда надо внутри \verb"Widget" хранить возможно не инициализированный \verb"Localizer". 
Это можно сделать либо через \verb"std::unique_ptr" либо через \verb"std::optional".
А так же нужно, чтобы \verb"Widget" корректно работал при не инициализированном локализаторе.
Так же можно сделать инициализацию и в конструкторе и по методу.

\item Singleton.
В данном случае мы просто берем локализатор из какого-то глобального хранилища данных.
Обычно такой подход  называют \verb"Service Locator".%
\footnote{Еще одно красивое название для дерьмового кода.}
\begin{cppcode}
class Widget {
public:
  ...
  void show() {
    ...
    auto& loc = ServiceLocator::instance()->WidgetLocalizator();
    title_ = loc->title();
    ...
  }
  ...
private:
  ...
};
\end{cppcode}

\item Observer-Observable.
Вместо неявной зависимости от \verb"Singleton" в виде \verb"ServiceLocator"-а, давайте соединим \verb"ServiceLocator" и \verb"Widget" с помощью push-уведомления через Observer-Observable паттерн.%
\footnote{Подробнее про него можно почитать в разделе~\ref{section::Observer}.}
То есть
\begin{cppcode}
class LocalizationModule {
public:
  void subscribe(Observer<Localizator>*);
private:
  Observable<Localizator> port_;
  Localizator loc_;
};

class Widget {
public:
  Observer<Localizator>* port();
  ...
  void show() {
    ...
    title_ = title();
    ...
  }
  ...
private:  
  std::string title() const {
    if (!port_.has_data())  // Need to support this case!
      return "";
    return port_->title();  // Take data from localizer received
  }
  ...
  Observer<Localizator> port_;
};
\end{cppcode}
И теперь нам лишь надо соединить два модуля вместе
\begin{cppcode}
LocalizationModule loc_module;
Widget widget;
loc_module.subscribe(widget.port());
\end{cppcode}
Все, теперь в этом месте у нас \verb"Widget" оповещается нужным локализатором и может его использовать.
\end{enumerate}
Давайте сделаем несколько замечаний по этим методам.
\begin{enumerate}
\item Методы $1$ и $2$ заставляют нас напрямую управлять объектом \verb"Widget" снаружи.
То есть если у нас сменился локализатор, то надо напрямую дернуть соответствующий метод \verb"Widget", чтобы изменения вступили в силу.

\item Метод $3$ страдает от неявной зависимости от единого \verb"ServiceLocator".
При таком подходе тяжело сменить какой-то сервис только для одного фиксированного объекта.
Если мы хотим поменять поведение какого-то сервиса внутри \verb"ServiceLocator", то это увидят все его пользователи.

\item Подход $4$ по сути является улучшенной версией подхода $3$, когда мы получаем данные не из глобального объекта, а из специального модуля, к которому мы явно присоединились.
Так же огромным достоинством этого подхода является push-оповещение.
То есть если пользователь меняет язык приложения, то внутри \verb"LocalizationModule" происходят необходимые изменения локализаторов и все подписчики мгновенно оповещаются об этом.

\item В подходе $3$ так же можно реализовать смену как самого \verb"ServiceLocator", так и отдельных его сервисов.
Однако, при такой замене НЕ происходит уведомление пользующихся им объектов.
А это значит, что придется ходить и руками дергать у всех пользователей \verb"ServiceLocator"-а служебные методы пересчета своего состояния.
Кроме того, проблема обостряется тем, что у вас нет никакого способа понять, а кто именно и где пользуется \verb"ServiceLocator"-ом, кроме как пройтись по всей кодовой базе.
\end{enumerate}
Я бы никогда не пользовался случаем $3$ за исключением ситуаций описанных в разделе~\ref{section::Global}, а точнее ситуациях аналогичных тем, что описаны в разделах~\ref{section::Logging} и~\ref{section::OpCounting}.
Например использование глобального Logger-а является хорошим примером для Singleton-а, потому что Logger -- это не часть кода, а это часть тестирующего кода (или собирающего статистику).
Методы $1$ и $2$ стоит использовать вместе и когда вы снаружи сами управляете объектом.
Метод $4$ я бы использовал как основной, только он позволяет автоматизировать своевременную реакцию на обновление зависимостей, позволяет менять зависимости налету и не содержит неявных связей между компонентами.
